\section{aemon: the agent as a product}\label{sec:app}
An agent that only exists in an evaluation harness is not useful to a manager. aemon runs the fine-tuned model (by default version~1 gathering the evidence and version~3 reasoning, with the
loop guards of Section~\ref{sec:results})
locally (FastAPI and React over a llama.cpp server) and is designed, after a review of how current research assistants
present long-running work, for readers who are not analysts.

\begin{itemize}
  \item \textbf{Answer first.} A finished investigation opens with the recommendation, a confidence label and score,
        the next step and what would change the answer, then four headline numbers, then the explanations compared,
        the test, the sources, and only then the full trace.
  \item \textbf{Numbers counted from the record.} Every chart is computed from the investigation's own record, never
        estimated: how likely each explanation was judged before and after new evidence (an arrow plot on the five-step
        scale), how many weighed sources support or contradict each explanation (diverging bars), the mix of
        independent sources and company claims, and the test's cost and duration against its budget.
  \item \textbf{A confidence score out of five.} One point for each of five plain checks: the leading explanation is
        judged likely or better; it is judged very likely; at least two independent sources support it; no independent
        source contradicts it; at least three sources were weighed for it. Five is High, three or four Moderate, two or
        fewer Low, and a draft written after a run stopped early is never more than Low. On the final test's 28 finished
        runs the score was 2 in 24 runs, 3 in two and 4 in two: the agent never judges its favourite very likely, and it
        weighs few independent sources, which is exactly what the iteration targets.
  \item \textbf{Follow-up questions grounded in the investigation.} A follow-up is answered by the base model (adapters
        off) from the investigation's record and source excerpts only, under a schema whose citations can name only
        sources in the record; ``Re-investigate with this'' reruns the investigation with the new information, files
        or links instead.
  \item \textbf{Steps a reader can follow.} The trace is shown as the seven steps a manager would recognise (understand
        the question, gather sources, list explanations, weigh the evidence, stress-test the favourite, design a test,
        recommend), each expandable to its actions and each action to exactly what was read, quoted or decided.
  \item \textbf{Investigations that do not die.} Runs keep going when the page is closed; every turn is saved, and
        after a restart unfinished runs resume by replaying their saved turns, each of which must reproduce its saved
        result. A run that stops early can be continued from where it stopped, and if it ends without a recommendation
        aemon writes a labelled draft from what it found.
  \item \textbf{Memory, projects, attachments and web search.} Notes about the person, the company or a project are
        given to the agent as context, never as evidence; attached files and links become quotable sources labelled as
        the person's own; an optional live web search saves every page it reads and labels it unchecked; a mind map
        shows every investigation; every result downloads as a PDF report.
\end{itemize}
None of these features is part of the evaluated setup; the reported results apply to the frozen configuration.

\section{Limitations}\label{sec:limits}
\begin{itemize}
  \item \textbf{Generalisation.} Four test companies, all business-to-business manufacturers asked the same family of
        question, support only a limited claim. The second test set addresses this partly; other industries and other
        kinds of strategic question remain untested.
  \item \textbf{Replications.} Three attempts per case at temperature 0 are close to repeats; the effective sample is
        nearer 16 cases than 48 runs, which is why the intervals resample whole cases. The cases are themselves 4
        companies with 4 variants of one question, identical until the what-if arrives; resampling whole companies
        widens the final-test interval to $[0.63, 1.05]$ (bootstrap) or $[0.44, 1.23]$ (cluster-robust $t$, 3 degrees of
        freedom), and test set 2's to $[0.29, 0.75]$ or $[0.25, 0.78]$. With 4 companies any interval is rough; all 16
        cases and all 4 companies are higher.
  \item \textbf{Judges.} Both judges are language models. They were qualified and agree closely on the order of runs,
        but not on the level (on the same fine-tuned runs, problem framing 1.62 from Claude and 2.90 from Codex), so the
        comparisons rest on the order and absolute scores mean little. They may share biases, and no strategy expert has
        scored the runs.
  \item \textbf{Absolute quality.} 1.48 of 3 is a large relative gain but a modest absolute score; 20 of 48 runs ended
        before a recommendation, and the recommendation, revision, experiment design and evidence judgement remain the
        weakest dimensions.
  \item \textbf{Comparators.} The comparison is with the same base model under its best prompt, which reads the full
        transcript: 25 of its 48 runs stop at the context limit, and a stage never reached scores 0. Comparing each pair
        only on the dimensions both runs were scored on gives $+0.92$ [0.83, 1.01] (test set 2: $+0.90$ [0.81, 0.99]), so
        the base model is weaker even on the stages it reaches (\texttt{reports/judge/robustness.json}). Neither the teacher nor
        a closed model was run on the final test for quality; Claude appears only in the cost comparison.
  \item \textbf{Beliefs and confidence.} Beliefs are the agent's own verbal scale, not calibrated probabilities, and the
        belief check tests direction only.
  \item \textbf{Sources and outcomes.} The dossier is fixed and curated, with short excerpts; what-if results are
        written by the case authors; no recommendation was checked against what a company later did, and no proposed
        test was run. The results measure the quality of the investigation, not whether its advice would pay off.
  \item \textbf{Training.} One training run and one seed; variance across runs is not measured. The training data has a
        known defect, the 240-character cut of test designs, and no example of declining a test.
  \item \textbf{Controls.} The scripted negative controls could not run on the held-out cases and are reported for
        development only.
  \item \textbf{Teacher errors in the training data.} The rows are valid, accepted calls, but their content carries
        the teacher's errors (Section~\ref{sec:negative}): invented numbers, wrong evidence kinds, overstated support,
        and, in historical training cases, sources published after the case date. The model reproduces the first two.
  \item \textbf{Our own instruments.} The mechanical experiment check is lenient: on development it passes 9 tests, 8
        of which fail stricter design screens. The evidence-kind annotation matches a blind relabelling on 75\% of records, 83\% counting defensible
        alternatives. Both are reported
        as they are, and the frozen numbers use them unchanged.
  \item \textbf{Mechanical validity.} Few runs pass every check: 1 of 48 on the final test and 4 of 72 on the second.
        The quality gains are measured by blind judges, not by mechanical validity.
\end{itemize}

\section{Conclusion}\label{sec:conclusion}
A small open model can learn a meaningful part of a disciplined strategy procedure from a larger model's tool-using
investigations, provided the environment makes the procedure executable and the data are filtered for exactness. On a
frozen test it framed problems better, built genuinely distinct causal theories, gathered and weighed quoted evidence and
finished far more investigations than its base model, at a fraction of the cost of a closed model. It did not learn what
its teacher rarely showed it: declining a test that is not worth its cost, responding to directional new evidence, and
writing complete decision rules, and it inherited the teacher's habits of unsupported numbers and loose evidence
labels. Neither runtime aids nor corrective examples from the same teacher repaired these, which places the remaining
gap in the quality of the supervision. A better teacher for evidence judgement is the next step. Equally important is how the result was obtained: an environment that refuses rather than trusts,
training data bound to its provenance, judges qualified before use, and a test frozen and run once.

\section*{Team and acknowledgements}
Hatem Saadallah and Nour Jennane, Master's students in artificial intelligence and AI researchers (team Two-eyed
Ravens). Built with AgentTune and AuditKIT (Lexsi Labs), Tiny Aya and Command A+ (Cohere), PEFT, TRL and llama.cpp.
Claude (Anthropic) and Codex (OpenAI) served as blind judges and as coding assistants; no Claude output was used as
training data.
